\chapter{Use Cases}
This chapter presents practical use cases in which sensor systems are applied to collect and analyze real-world data. Different application domains demonstrate how sensor technologies are integrated into larger systems to monitor physical conditions, human activity, and environmental parameters.

\section{Mobile Movement Capturing}
Wearable's/Fitness Analytics Devices incorporate a variety of sensors to keep track of several parameters of the person wearing one of those devices. The measurements recorded are then put into various algorithms to determine a more accurate state.

\boldtitle{IMU}
Movement Capturing Devices use IMUs to determine the position and velocity of the user. This is especially useful for tracking steps, calculating the path, the individual moved along and even determining the inclination of the user, when walking.

\noindent\cite{LRWeb:47}

\boldtitle{Altimeter}
Some Fitness Analytics Devices incorporate a component for sensing the altitude. This is mostly used in mountain climbing and it works by measuring the pressure, the sensor is experiencing. This lets the analytics describe the path more accurately since it can sense elevation and gradient in a path. 

\noindent\cite{LRWeb:47}

\boldtitle{Temperature Sensors \& Pulse Sensors}
Belong to the category of Biomedical sensors.

The temperature sensor monitors body temperature and can trigger warnings if it rises too high, using measurement principles similar to conventional temperature sensors.

\noindent\cite{LRWeb:47}

Pulse sensors emit green light into the skin and changes in blood flow alter how much light is reflected back.

\noindent\cite{LRWeb:48}

\section{Environmental monitoring}
Data centers accommodate large amounts of expensive hardware-infrastructure. Therefore, it is economically wise to ensure that those components work in the best condition and are maintained correctly. The correct temperature and humidity are key parts for excellent working conditions. To optimize these conditions, there are several steps to achieve this.

\boldtitle{Collecting and analyzing real-time data}
Sensors continuously collect environmental sensor data in regular intervals. These sensors are strategically placed in even distances or in regions where more activity is happening. After the data is collected, it is analyzed to detect trends and anomalies. The data is compared to previously defined thresholds and if they exceed or fall below those thresholds, an alarm is triggered.

\boldtitle{Taking actions on alarm}
When an alarm is emitted, a HVAC\footnote{Heating, Ventilation and Air Conditioning}-system are triggered to achieve optimal conditions.

\boldtitle{Ensuring optimal conditions}
To make sure that optimal conditions are always met, additional redundant sensors can also be placed, to backup possible sensor outages.

\noindent\cite{LRWeb:56}

\section{Smart Infrastructure (Smart Home)}
\begin{quote}
    Smart infrastructure is the use of digital technologies to improve the efficiency, sustainability, and resilience of physical infrastructure. It involves collecting and analyzing data from sensors and other devices to better understand how infrastructure is being used and where it is under strain.
\end{quote}

\noindent\cite{LRWeb:45}

The smart infrastructure is usually paired with a mobile app to control the state of the infrastructure. It may also be that automated functionality is executed on events caused by the sensors, e.g.\ call authorities when a presence sensor detects unwanted movement inside the home.

Smart Home sensors perform a variety of functions such as monitoring doors, windows, movement, temperature, and humidity. Common examples include contact sensors for opening and closing states, motion and presence sensors for detecting occupancy and climate sensors for environmental conditions. These sensor types are described in consumer-oriented technical guides on Smart Home technologies such as the Euronics Trendblog overview of Smart Home sensors

\boldtitle{Contact Sensor}
A Contact Sensor consists of two magnetic pieces that are usually placed on windows or doors. When the distance between those two components gets too large, an event occurs that notifies the Smart Home about the change. This sensor can be used to stop heating when the windows are opened to ventilate the home. Another use case might be to send a notification when an intruder enters the home.

\boldtitle{Motion Sensor}
A Motion Sensor sends out infrared-light in a broad angle. This lets the sensor detect changes in temperature if somebody walks past it. This sensor is usually used to turn on the light if a person enters the room.

A more precise version of the Motion Sensor is the Presence Sensor. This component operates by sending millimeter-waves which can detect the smallest movements. This sensor is usually more expensive but much more accurate.

\boldtitle{Humidity Sensors}
A poor proportion of humidity in a room causes trouble for humans (described in~\ref{title:humidity}). Humidity Sensors notify Smart Homes, when it is time to ventilate.

\noindent\cite{LRWeb:46}